\SourceLeaf{19}
into rectangular elements. Such an element is enclosed by the points with coordinates
\[
\begin{array}{ll}
p,&q\\
p+\mathrm{d}p,&q\\
p,&q+\mathrm{d}q\\
p+\mathrm{d}p,&q+\mathrm{d}q
\end{array}
\]
Its area will be $=\mathrm{d}p.\mathrm{d}q$, and it corresponds to a parallelogram-shaped element on the surface of the body, enclosed by four points whose respective coordinates will be
\[
\begin{array}{rlll}
\mathrm{I.}&x,&y,&z\\
\mathrm{II.}&x+\lambda\mathrm{d}p,&y+\mu\mathrm{d}p,&z+\nu\mathrm{d}p\\
\mathrm{III.}&x+\lambda'\mathrm{d}q,&y+\mu'\mathrm{d}q,&z+\nu'\mathrm{d}q\\
\mathrm{IV.}&x+\lambda\mathrm{d}p+\lambda'\mathrm{d}q,&y+\mu\mathrm{d}p+\mu'\mathrm{d}q,&z+\nu\mathrm{d}p+\nu'\mathrm{d}q
\end{array}
\]
if we assume that
\[
\begin{aligned}
\mathrm{d}x&=\lambda\mathrm{d}p+\lambda'\mathrm{d}q\\
\mathrm{d}y&=\mu\mathrm{d}p+\mu'\mathrm{d}q\\
\mathrm{d}z&=\nu\mathrm{d}p+\nu'\mathrm{d}q
\end{aligned}
\]
The projections of this area, which we set $=\mathrm{d}s$, onto the three planes normal to the coordinate axes $x,y,z$ are easily found to be, respectively, $=$
\[
\begin{aligned}
&\pm(\mu\nu'-\nu\mu')\mathrm{d}p.\mathrm{d}q\\
&\pm(\nu\lambda'-\lambda\nu')\mathrm{d}p.\mathrm{d}q\\
&\pm(\lambda\mu'-\mu\lambda')\mathrm{d}p.\mathrm{d}q
\end{aligned}
\]
so by a well-known theorem the area of the element itself will be
\[
=\mathrm{d}p.\mathrm{d}q.\BareRoot{((\mu\nu'-\nu\mu')^2+(\nu\lambda'-\lambda\nu')^2+(\lambda\mu'-\mu\lambda')^2)}
\]
\par\endgroup
\par
It follows that each integral given in our six theorems can be reduced to the form $\int S\mathrm{d}p.\mathrm{d}q$, where $S$ is either explicitly or implicitly a function of the two indeterminates $p,q$. The integration is to be extended either over the whole infinite plane or over the portion of it that represents, as it were, the entire surface of our body. The integration itself will be carried out by one device or another, for which no general rules can be given.
